\documentclass{article}

% Project-specific macros — exercise the extractor on the patterns we care about.
\usepackage{amsmath}
\usepackage{amssymb}
\usepackage{amsthm}
\usepackage{mathtools}
\usepackage{cleveref}

\newcommand{\R}{\mathbb{R}}
\newcommand{\N}{\mathbb{N}}
\newcommand{\norm}[1]{\left\lVert #1 \right\rVert}
\newcommand{\inner}[2]{\langle #1, #2 \rangle}
\DeclareMathOperator{\Tr}{Tr}
\DeclareMathOperator*{\argmin}{arg\,min}
\providecommand{\eps}{\varepsilon}
\def\Z{\mathbb{Z}}

\usepackage{mathpreview}

\title{Convergence of $\sum 1/n^s$}
\author{TexViewer Demo}

\begin{document}
\maketitle

\section{Introduction}
\label{sec:intro}
We study the convergence of the series $\sum_{n=1}^{\infty} 1/n^s$ over the
real numbers $s \in \R$. The result is classical, but the proof exercises a
handful of standard inequalities that we collect in \cref{sec:lemmas}.

\section{Main result}
\label{sec:main}

\begin{theorem}[role=main]{Convergence of the zeta-like series}
\label{thm:main}
For every $s \in \R$, the series
\[
  \zeta(s) \;=\; \sum_{n=1}^{\infty} \frac{1}{n^s}
\]
converges if and only if $s > 1$.
\end{theorem}

\begin{proof}[of \cref{thm:main}]
By \cref{lem:harmonic} the partial sums $H_N$ diverge for $s = 1$, so the
series cannot converge there. For $s > 1$ the integral test
(\cref{lem:integral-test}) gives a finite bound. For $s < 1$ comparison with
the harmonic series settles it. \cite{Rudin1976}
\end{proof}

\section{Supporting lemmas}
\label{sec:lemmas}

\begin{lemma}[role=supporting]{Harmonic divergence}
\label{lem:harmonic}
For every $N \ge 1$, $H_N := \sum_{n=1}^{N} 1/n \ge \log(N+1)$. In
particular $H_N \to \infty$.
\end{lemma}

\begin{proof}
The function $1/x$ is decreasing on $(0,\infty)$, so for each integer
$n \ge 1$,
\[
  \frac{1}{n} \;\ge\; \int_{n}^{n+1} \frac{dx}{x}.
\]
Summing $n = 1, \dots, N$ yields $H_N \ge \int_1^{N+1} \frac{dx}{x} = \log(N+1)$. The
right-hand side diverges as $N \to \infty$.
\end{proof}

\begin{lemma}[role=standard]{Integral test}
\label{lem:integral-test}
Let $f \colon [1, \infty) \to (0, \infty)$ be continuous and decreasing.
Then $\sum_{n=1}^{\infty} f(n)$ converges iff $\int_1^{\infty} f(x)\,dx$
converges.
\end{lemma}

\begin{proof}
Standard; see any first-year analysis text.
\end{proof}

\begin{theorem}[role=omitted]{Cauchy condensation}
\label{thm:cauchy-cond}
Let $(a_n)$ be a non-increasing sequence of non-negative reals. Then
$\sum a_n$ converges iff $\sum 2^k a_{2^k}$ converges.
\omitref{Rudin, \emph{Principles of Mathematical Analysis}, §3.27}
\end{theorem}

\section{An assorted identity}

A toy display to exercise custom macros:
\begin{equation}
\label{eq:innerprod}
  \norm{x}^2 \;=\; \inner{x}{x}, \qquad \Tr(AB) = \Tr(BA), \qquad
  \argmin_{x \in \R^n} \norm{Ax - b}^2 = (A^\top A)^{-1} A^\top b.
\end{equation}

We will refer back to \eqref{eq:innerprod} in passing, and use $\eps > 0$ as
a convergence tolerance. Note that $\Z \subset \R$.

\end{document}
